\input{../tex-inputs/latex-leseansicht-vorspann}

\section[Rainer Maria Rilke: Widmungsexemplar Das tägliche Leben für Arthur Schnitzler, {[}18. 1. 1902{]}]{L04545 Rainer Maria Rilke: Widmungsexemplar Das tägliche Leben für Arthur
               Schnitzler, {[}18. 1. 1902{]}}
\nopagebreak\mylabel{L04545v}
\rehead{ }\normalsize\beginnumbering\briefempfaengerindex{Schnitzler, Arthur@\textsc{Schnitzler, Arthur}!zzzRilke, Rainer Maria@\emph{von Rainer Maria Rilke}!1902-01-182@{{[}18. 1. 1902{]}}|(be}
\toendnotes[C]{\smallbreak\pagebreak[2]}
\correspDesc{Versand  durch Rainer Maria Rilke am [18. 1. 1902] in Westerwede
\newline{}Erhalt  durch Arthur Schnitzler im Zeitraum [19. 1. 1902 – 23. 1. 1902?] in Wien}\toendnotes[C]{\smallbreak}
\buchAlsQuelle{\emph{Rainer Maria Rilke und Arthur Schnitzler. Ihr Briefwechsel.}. Mit Anmerkungen herausgegeben von Heinrich Schnitzler In: \emph{Wort und Wahrheit}, Jg. 13, Nr. 4, April 1958, S. 282–298, hier S. 298.}\toendnotes[C]{\smallbreak}
\pstart
           \noindent{}{\pb}\pwindex{Rilke, Rainer Maria 4.\,12.\,1875 Prag – 29.\,12.\,1926 Montreux@\textsc{Rilke, Rainer Maria} (4.\,12.\,1875 Prag – 29.\,12.\,1926 Montreux)!tägliche Leben. Drama in zwei Akten@\strich\emph{Das tägliche Leben. Drama in zwei Akten}|pw}\label{K_L04539-1v}\edtext{An Arthur Schnitzler}{\lemma{\textnormal{\emph{An Arthur Schnitzler}}}\Cendnote{\textnormal{Zur Datierung siehe XXXX Auszeichnungsfehler: Dokument L04543 nicht gefunden. Das Widmungsexemplar befand sich 1958 im Besitz
                  von Heinrich Schnitzler\pwindex{Schnitzler, Heinrich 9.\,8.\,1902 Hinterbrühl – 12.\,7.\,1982 Wien@\textsc{Schnitzler, Heinrich} (9.\,8.\,1902 Hinterbrühl – 12.\,7.\,1982 Wien)|pwk}. Der Verbleib ist
                  ungeklärt.}}}\label{K_L04539-1}, Rainer Maria Rilke in herzlicher Verehrung\pend
           \selectlanguage{ngerman}\endnumbering\briefempfaengerindex{Schnitzler, Arthur@\textsc{Schnitzler, Arthur}!zzzRilke, Rainer Maria@\emph{von Rainer Maria Rilke}!1902-01-182@{{[}18. 1. 1902{]}}|)be}\mylabel{L04545h}
\begin{anhang}
\end{anhang}\newcommand{\dateiname}{L04545}\newcommand{\titel}{Rainer Maria Rilke: Widmungsexemplar Das tägliche Leben für Arthur Schnitzler, [18. 1. 1902]}\newcommand{\editorInnen}{Herausgegeben von Jahnke, SelmaMüller, Martin Anton}\input{../tex-inputs/latex-leseansicht-abspann}
